% GENERATED FILE. Edit canonical lessons, then run npm run generate:books.
% canonical-source-hash: fnv1a64:bc95fec7d51645dd
% canonical-lessons: KA-C412-apaya, KA-C412-munneccarike, KA-C412-sambhavaniya, KA-C412-sankasta, KA-C412-nivarisu

\chapter{When Danger Comes}
\label{ch:ka-when-danger-comes}
\hlchaptermodality{412}

\begin{chapteropening}
\textbf{By the end of this chapter:} \emph{I can say a danger, a precaution, probable, hardship and to ward off.}

Complete the last lesson of chapter 412: I can say a danger, a precaution, probable, hardship and to ward off.
\end{chapteropening}

\section[\texorpdfstring{\kn{ಅಪಾಯ} (apāya)}{ಅಪಾಯ (apāya)}]{\kn{ಅಪಾಯ} (apāya) — a danger; a risk}

\label{lesson:KA-C412-apaya}



\begin{quote}
Before the new one: say the Kannada for to fulfil, then the Kannada for an obligation.
\end{quote}

\subsection*{You'll want to know: \kn{ಅಪಾಯ}}
\textbf{\kn{ಅಪಾಯ}} — \emph{apāya} — \textquotedblleft{}a danger; a risk\textquotedblright{}.

\begin{quote}
\kn{ಮಳೆಗಾಲದಲ್ಲಿ} \kn{ನದಿಯಲ್ಲಿ} \kn{ಈಜುವುದು} \kn{ಅಪಾಯ}.
\end{quote}

\begin{quote}
\emph{maḷegāladalli nadiyalli ījuvudu apāya.}
\end{quote}

\begin{quote}
Swimming in the river in the rainy season is dangerous.
\end{quote}

\kn{ಅಪಾಯದಲ್ಲಿ} \kn{ಇರು}, to be in danger.

\begin{cousinweb}
\textbf{\kn{ಅಪಾಯ}} (\emph{apāya}) is borrowed from Sanskrit \emph{apāya}, \textquotedblleft{}going away; loss, harm\textquotedblright{} (CDIAL 481).
\end{cousinweb}

\begin{grammarlens}[title={What you've built}]
One more word for saying what happens if.
\end{grammarlens}

\subsection*{Your turn}
\begin{itemize}
\raggedright
  \item \emph{Say it:} \emph{apāya}
  \item \emph{Say it:} \emph{apāya}, once more
  \item \emph{Say it:} \emph{maḷegāladalli nadiyalli ījuvudu apāya.} — Swimming in the river in the rainy season is dangerous.
  \item \emph{Recall:} say the Kannada for to fulfil, then the Kannada for an obligation, then say \emph{apāya} again
\end{itemize}

\begin{tcolorbox}[breakable,colback=teal!4,colframe=teal!35!black,title={Before you move on}]
What does \kn{ಅಪಾಯ} mean? (\textquotedblleft{}A danger; a risk\textquotedblright{}.) Say it once more.
\end{tcolorbox}
\section[\texorpdfstring{\kn{ಮುನ್ನೆಚ್ಚರಿಕೆ} (munneccarike)}{ಮುನ್ನೆಚ್ಚರಿಕೆ (munneccarike)}]{\kn{ಮುನ್ನೆಚ್ಚರಿಕೆ} (munneccarike) — a precaution}

\label{lesson:KA-C412-munneccarike}



\begin{quote}
Before the new one: say the Kannada for an obligation, then the Kannada for a danger.
\end{quote}

\subsection*{You'll want to know: \kn{ಮುನ್ನೆಚ್ಚರಿಕೆ}}
\textbf{\kn{ಮುನ್ನೆಚ್ಚರಿಕೆ}} — \emph{munneccarike} — \textquotedblleft{}a precaution\textquotedblright{}.

\begin{quote}
\kn{ಮುನ್ನೆಚ್ಚರಿಕೆಯಾಗಿ} \kn{ನಾನು} \kn{ಕೊಡೆ} \kn{ತೆಗೆದುಕೊಂಡು} \kn{ಹೋದೆ}.
\end{quote}

\begin{quote}
\emph{munneccarikeyāgi nānu koḍe tegedukoṇḍu hōde.}
\end{quote}

\begin{quote}
I took an umbrella as a precaution.
\end{quote}

\kn{ಮುನ್} (before) + \kn{ಎಚ್ಚರಿಕೆ} (warning).

\begin{cousinweb}
\textbf{\kn{ಮುನ್ನೆಚ್ಚರಿಕೆ}} (\emph{munneccarike}) is native Dravidian, built from \emph{mun}, \textquotedblleft{}before\textquotedblright{} (DEDR 5020), and \emph{eccaṟike}, \textquotedblleft{}a warning, watchfulness\textquotedblright{} (DEDR 851).
\end{cousinweb}

\begin{grammarlens}[title={What you've built}]
One more word for saying what happens if.
\end{grammarlens}

\subsection*{Your turn}
\begin{itemize}
\raggedright
  \item \emph{Say it:} \emph{munneccarike}
  \item \emph{Say it:} \emph{munneccarike}, once more
  \item \emph{Say it:} \emph{munneccarikeyāgi nānu koḍe tegedukoṇḍu hōde.} — I took an umbrella as a precaution.
  \item \emph{Recall:} say the Kannada for an obligation, then the Kannada for a danger, then say \emph{munneccarike} again
\end{itemize}

\begin{tcolorbox}[breakable,colback=teal!4,colframe=teal!35!black,title={Before you move on}]
What does \kn{ಮುನ್ನೆಚ್ಚರಿಕೆ} mean? (\textquotedblleft{}A precaution\textquotedblright{}.) Say it once more.
\end{tcolorbox}
\section[\texorpdfstring{\kn{ಸಂಭವನೀಯ} (sambhavanīya)}{ಸಂಭವನೀಯ (sambhavanīya)}]{\kn{ಸಂಭವನೀಯ} (sambhavanīya) — probable; likely}

\label{lesson:KA-C412-sambhavaniya}



\begin{quote}
Before the new one: say the Kannada for a danger, then the Kannada for a precaution.
\end{quote}

\subsection*{You'll want to know: \kn{ಸಂಭವನೀಯ}}
\textbf{\kn{ಸಂಭವನೀಯ}} — \emph{sambhavanīya} — \textquotedblleft{}probable; likely\textquotedblright{}.

\begin{quote}
\kn{ಮಳೆ} \kn{ಬಂದರೆ} \kn{ಪ್ರವಾಹ} \kn{ಬರುವುದು} \kn{ಸಂಭವನೀಯ}.
\end{quote}

\begin{quote}
\emph{maḷe bandare pravāha baruvudu sambhavanīya.}
\end{quote}

\begin{quote}
If it rains, a flood is likely.
\end{quote}

\kn{ಸಂಭವ}, a possibility or an event.

\begin{cousinweb}
\textbf{\kn{ಸಂಭವನೀಯ}} (\emph{sambhavanīya}) is built on Sanskrit \emph{saṃbhava}, \textquotedblleft{}coming together, coming into being\textquotedblright{} (Monier-Williams), with -\emph{anīya}, \textquotedblleft{}that may be\textquotedblright{}.
\end{cousinweb}

\begin{grammarlens}[title={What you've built}]
One more word for saying what happens if.
\end{grammarlens}

\subsection*{Your turn}
\begin{itemize}
\raggedright
  \item \emph{Say it:} \emph{sambhavanīya}
  \item \emph{Say it:} \emph{sambhavanīya}, once more
  \item \emph{Say it:} \emph{maḷe bandare pravāha baruvudu sambhavanīya.} — If it rains, a flood is likely.
  \item \emph{Recall:} say the Kannada for a danger, then the Kannada for a precaution, then say \emph{sambhavanīya} again
\end{itemize}

\begin{tcolorbox}[breakable,colback=teal!4,colframe=teal!35!black,title={Before you move on}]
What does \kn{ಸಂಭವನೀಯ} mean? (\textquotedblleft{}Probable; likely\textquotedblright{}.) Say it once more.
\end{tcolorbox}
\section[\texorpdfstring{\kn{ಸಂಕಷ್ಟ} (saṅkaṣṭa)}{ಸಂಕಷ್ಟ (saṅkaṣṭa)}]{\kn{ಸಂಕಷ್ಟ} (saṅkaṣṭa) — hardship; distress}

\label{lesson:KA-C412-sankasta}



\begin{quote}
Before the new one: say the Kannada for a precaution, then the Kannada for probable.
\end{quote}

\subsection*{You'll want to know: \kn{ಸಂಕಷ್ಟ}}
\textbf{\kn{ಸಂಕಷ್ಟ}} — \emph{saṅkaṣṭa} — \textquotedblleft{}hardship; distress\textquotedblright{}.

\begin{quote}
\kn{ಪ್ರವಾಹದ} \kn{ನಂತರ} \kn{ಹಳ್ಳಿಯ} \kn{ಜನರು} \kn{ಸಂಕಷ್ಟದಲ್ಲಿ} \kn{ಇದ್ದರು}.
\end{quote}

\begin{quote}
\emph{pravāhada nantara haḷḷiya janaru saṅkaṣṭadalli iddaru.}
\end{quote}

\begin{quote}
After the flood the villagers were in distress.
\end{quote}

\kn{ಕಷ್ಟ} with \kn{ಸಂ}-: a heavy trouble.

\begin{cousinweb}
\textbf{\kn{ಸಂಕಷ್ಟ}} (\emph{saṅkaṣṭa}) is borrowed from Sanskrit \emph{saṃkaṣṭa}, \textquotedblleft{}distress, trouble\textquotedblright{} (Monier-Williams), built on \emph{kaṣṭa}, \textquotedblleft{}hardship\textquotedblright{} (CDIAL 2978).
\end{cousinweb}

\begin{grammarlens}[title={What you've built}]
One more word for saying what happens if.
\end{grammarlens}

\subsection*{Your turn}
\begin{itemize}
\raggedright
  \item \emph{Say it:} \emph{saṅkaṣṭa}
  \item \emph{Say it:} \emph{saṅkaṣṭa}, once more
  \item \emph{Say it:} \emph{pravāhada nantara haḷḷiya janaru saṅkaṣṭadalli iddaru.} — After the flood the villagers were in distress.
  \item \emph{Recall:} say the Kannada for a precaution, then the Kannada for probable, then say \emph{saṅkaṣṭa} again
\end{itemize}

\begin{tcolorbox}[breakable,colback=teal!4,colframe=teal!35!black,title={Before you move on}]
What does \kn{ಸಂಕಷ್ಟ} mean? (\textquotedblleft{}Hardship; distress\textquotedblright{}.) Say it once more.
\end{tcolorbox}
\section[\texorpdfstring{\kn{ನಿವಾರಿಸು} (nivārisu)}{ನಿವಾರಿಸು (nivārisu)}]{\kn{ನಿವಾರಿಸು} (nivārisu) — to ward off; to relieve}

\label{lesson:KA-C412-nivarisu}



\begin{quote}
Before the new one: say the Kannada for probable, then the Kannada for hardship.
\end{quote}

\subsection*{You'll want to know: \kn{ನಿವಾರಿಸು}}
\textbf{\kn{ನಿವಾರಿಸು}} — \emph{nivārisu} — \textquotedblleft{}to ward off; to relieve\textquotedblright{}.

\begin{quote}
\kn{ವೈದ್ಯರು} \kn{ನನ್ನ} \kn{ನೋವನ್ನು} \kn{ನಿವಾರಿಸಿದರು}.
\end{quote}

\begin{quote}
\emph{vaidyaru nanna nōvannu nivārisidaru.}
\end{quote}

\begin{quote}
The doctor relieved my pain.
\end{quote}

Past (I …): \kn{ನಿವಾರಿಸಿದೆ}, \emph{nivāriside}.

What is removed takes -\kn{ಅನ್ನು}. \kn{ನಿವಾರಣೆ}, removal.

\begin{cousinweb}
\textbf{\kn{ನಿವಾರಿಸು}} (\emph{nivārisu}) is built on Sanskrit \emph{nivāraṇa}, \textquotedblleft{}keeping off, warding off\textquotedblright{} (CDIAL 7418), with the Kannada verb ending -\emph{isu}.
\end{cousinweb}

\begin{grammarlens}[title={What you've built}]
One more word for saying what happens if.
\end{grammarlens}

\subsection*{Your turn}
\begin{itemize}
\raggedright
  \item \emph{Say it:} \emph{nivārisu}
  \item \emph{Say it:} \emph{nivārisu}, once more
  \item \emph{Say it:} \emph{vaidyaru nanna nōvannu nivārisidaru.} — The doctor relieved my pain.
  \item \emph{Recall:} say the Kannada for probable, then the Kannada for hardship, then say \emph{nivārisu} again
\end{itemize}

\begin{tcolorbox}[breakable,colback=teal!4,colframe=teal!35!black,title={Before you move on}]
What does \kn{ನಿವಾರಿಸು} mean? (\textquotedblleft{}To ward off; to relieve\textquotedblright{}.) Say it once more.
\end{tcolorbox}
